\documentclass[12pt]{article}

\usepackage{amsmath}
\usepackage{amssymb}
\usepackage{graphicx}
\usepackage{geometry}
\usepackage{setspace}
\usepackage{hyperref}
\usepackage{iwona}
\usepackage{natbib}
\usepackage{booktabs}
\usepackage{array}
\usepackage{xcolor}

\geometry{a4paper, margin=1in}
\onehalfspacing

\hypersetup{
  colorlinks=true,
  linkcolor=black,
  citecolor=black,
  urlcolor=black
}

\title{Flow, Autonomy, and Organizational Heat: \\
A Physical Systems Analogy for Psychology, Decision Architecture,
and Organization Design}

\author{Rickard Hultgren}
\date{\today}

\begin{document}

\maketitle

\begin{abstract}
Organizational systems frequently exhibit patterns of friction, conflict, and systemic
overload that resist explanation through conventional accounts of motivation, leadership,
or communication. This paper develops a conceptual framework, grounded in hydraulic and
electrical circuit analogies, to model how autonomy, legitimacy, capability, risk, and
structural complexity interact to produce either smooth coordinated action or what we
term \textit{organizational heat}---the dissipative byproduct of action flow encountering
structural resistance. The framework proposes five core formal relations as modeling hypotheses: autonomy as the product of legitimacy and capability; legitimacy and capability as bounded quantities subject to reduction by insignificance and risk, respectively; action flow as a ratio of autonomy to complexity resistance; and organizational heat as proportional to the square of flow times resistance. It further distinguishes between legitimacy as durable standing within a decision domain and admissibility as the execution-time validity of a specific act under current conditions. These algebraic forms are the
authors' own formalizations by analogy to physical circuit theory and are not directly
derived from the cited literatures, which are invoked to motivate the conceptual
components rather than to validate the specific functional forms. The model is extended through Kirchhoff-like network constraints applied to decision architectures, through a distinction between coordination flow and irreversible commitment, and through analysis of how AI-augmented organizations alter the flow--resistance--heat balance while increasing the importance of commit-time admissibility checks.
We conclude with testable predictions that could validate or refute the framework's
structural assumptions, and with an explicit assessment of model limitations.
\end{abstract}

\tableofcontents
\newpage

\section{Introduction}

\subsection{The Problem of Organizational Instability}

Organizations routinely experience coordination failures, decision overload,
rework, and escalating interpersonal conflict even when their members are
individually competent and their objectives are clearly articulated. Standard
accounts of such failures frequently invoke motivational deficits
\citep{deci1985intrinsic}, leadership dysfunction \citep{yukl2012leadership},
or communication breakdowns \citep{weick1995sensemaking}. While each of these
accounts captures important phenomena, they do not adequately explain a class
of instabilities that arise structurally---that is, as a consequence of how
decision authority, information load, workflow routing, and interdependency are
arranged, rather than as a consequence of individual deficiency or
interpersonal failure.

Physical systems provide a potentially productive set of analogies for this
structural class of problem. In fluid and electrical networks, flow is
governed by driving potential and opposed by resistance; when flow encounters
resistance, energy is dissipated as heat \citep{feynman1963lectures}. This
paper proposes that analogous structural constraints may apply to the flow of
decisions and actions within organizations. The specific algebraic forms
through which we formalize this analogy---presented fully in
Sections~\ref{sec:autonomy}--\ref{sec:heat}---are modeling choices rather than
empirical findings; they are introduced because they generate tractable,
falsifiable predictions, not because they are established by the literatures
cited for motivation.

A further refinement is necessary. In complex organizations, the first visible
symptoms of failure often appear in execution---for example, on the shop
floor, in service delivery, in customer outcomes, or in compliance breaches.
Yet execution commonly reveals the symptom rather than the point of origin.
Likewise, data inconsistencies, reporting anomalies, or KPI conflicts may
appear to be the first breakdown, while in fact they often reflect a prior
distortion in how the underlying work was decomposed, routed, sequenced,
escalated, or reconciled. The present framework therefore gives analytic
priority to the \emph{workflow or orchestration layer}: the intermediate
structure through which strategic and managerial intent is translated into
tasks, handoffs, system actions, escalation paths, and coordinated execution.

This distinction matters because organizations do not act directly on strategy.
They act through orchestrated workflows. If the workflow or orchestration layer
encodes a slightly altered version of the original decision, then downstream
actors, systems, and data structures may all behave in ways that are locally
consistent while globally misaligned with the organization's actual intent.
The present paper therefore argues that many coordination failures should be
understood not simply as failures of communication or authority, but as
failures of \emph{intent preservation through orchestration}.

A second complication arises in distributed and partially automated systems:
even when decisions are initially well-founded and even when orchestration
remains coherent, the conditions that made a specific action valid may change
before the action is actually committed. This implies that organizations face
not one but two distinct control problems. First, they must preserve coherent
intent as decisions are translated into workflows across humans, teams,
software systems, and AI agents. Second, they must determine whether a
specific act remains admissible at the moment it becomes irreversible. The
present paper argues that both problems must be modeled explicitly and that
they occur at different structural layers.

\subsection{Literature Positioning}

The framework draws on five bodies of established theory to motivate its
conceptual components. We note explicitly where the framework is aligned with
existing evidence and where it extends beyond what those sources directly
support.

Across these literatures, authority, control, and validity are often discussed
without sharply distinguishing durable standing from execution-time
authorization, and workflow structure is often treated as an implementation
detail rather than as a primary site of strategic distortion. The present
framework departs from this tendency in two ways. First, it distinguishes
durable standing legitimacy from execution-time admissibility. Second, it
treats workflow or orchestration architecture as the principal intermediate
layer at which intent is either preserved or lost before downstream data and
execution outcomes make the loss visible.

\paragraph{Self-Determination Theory (SDT).}
Deci and Ryan's self-determination theory \citep{deci2000self, ryan2000self}
identifies autonomy, competence, and relatedness as fundamental psychological
needs whose satisfaction reliably predicts intrinsic motivation and well-being
across a wide range of empirical studies. This evidence base motivates our
choice of autonomy as the primary driving variable in the model. The present
framework departs from SDT in treating autonomy as a structural quantity with
a specific algebraic form (the product of legitimacy and capability), rather
than as a psychological need whose satisfaction falls on a continuous scale.
That algebraic form is a modeling assumption, not something SDT or its
associated empirical literature directly implies.

\paragraph{Cognitive Load Theory (CLT).}
Sweller's cognitive load theory \citep{sweller1988cognitive,
sweller1994cognitive} establishes that working memory is a finite resource and
that performance degrades when total processing demand---across its intrinsic,
extraneous, and germane components---exceeds available capacity. This finding
supports the general principle that complexity imposes a cost on effective
action, which motivates the resistance construct in the present model. The
specific claim that workflow orchestration and organizational complexity
function jointly as resistance in a ratio-form flow equation is our own
formalization, not a conclusion of CLT.

\paragraph{Systems Theory and Cybernetics.}
General systems theory \citep{bertalanffy1968general} and cybernetics
\citep{ashby1956introduction, wiener1948cybernetics} model organizations as
open systems regulated by feedback, characterized by flows of energy,
information, and material. This tradition motivates the use of flow and
network concepts in describing organizations. The application of
Kirchhoff-like conservation constraints to decision networks---developed in
Section~\ref{sec:kirchhoff}---is our extension of this tradition to the domain
of authority and coordination flows; it is a conceptual analogy, not an
established result in organizational systems theory.

\paragraph{Control Theory and Normal Accidents.}
Perrow's normal accident theory \citep{perrow1984normal} and Beer's viable
systems model \citep{beer1972brain} both establish that organizational failure
is systematically produced by structural mismatches between complexity and
regulatory capacity. This provides empirical and theoretical grounding for the
claim that structural properties, rather than individual behavior alone,
determine system stability. The present framework extends this claim by
proposing that one important locus of early instability is the workflow or
orchestration layer, where regulatory intent is first transformed into
operational sequence, routing, and handoff logic. The further claim that such
instability manifests as a quadratic heat function remains a model prediction
that requires empirical testing against the failure-mode patterns that Perrow
and Beer document.

\paragraph{Coordination Costs and Organization Design.}
Transaction cost economics \citep{williamson1981economics} and coordination
theory \citep{malone1994interdisciplinary} establish that organizational design
choices have systematic effects on the costs of integrating distributed action.
\citet{puranam2012coordination} demonstrate that coordination failure is
predictably linked to specific structural properties of task interdependence.
The present model builds on this tradition by treating coordination cost not
only as a static burden but as a dynamic resistance parameter that interacts
with autonomy to produce either stable throughput or dissipative heat. It
further argues that in high-complexity systems this burden is mediated
centrally by orchestration architecture: the structural layer where decisions
are decomposed into executable paths across roles and systems.

\subsection{Overview of the Framework}

The model proposes the following formal relations and architectural
distinctions as a coherent set of modeling hypotheses:

\begin{enumerate}
  \item Autonomy as the product of legitimacy and capability
  \item Legitimacy as a bounded quantity reduced by insignificance
  \item Capability as a bounded quantity reduced by risk
  \item Action flow as autonomy divided by complexity resistance
  \item Organizational heat as the square of flow multiplied by resistance
  \item Kirchhoff-like conservation constraints on decision networks
  \item Decision architecture as the structural allocation of standing authority
  \item Orchestration architecture as the structural layer through which intent is decomposed, routed, sequenced, escalated, and either preserved or distorted
  \item Commit architecture as the determinant of whether action remains admissible at the moment of irreversible execution
\end{enumerate}

These relations do not all operate at the same layer. The first six concern
the generation, routing, and dissipation of coordinated action through an
organizational network. The seventh concerns who possesses recognized standing
to decide within a domain. The eighth concerns how those decisions are
translated into coordinated workflows across human and technical components.
The ninth concerns a distinct but related problem: whether a specific action
remains admissible at the point of commitment. Their value as a theoretical
contribution depends not on the precision of the analogy to physics but on
whether they generate predictions that can be tested against organizational
data. Section~\ref{sec:predictions} develops these predictions explicitly.

Figure~\ref{fig:model} presents the integrated model.

\begin{figure}[h]
\centering
\includegraphics[width=0.9\textwidth]{graph.pdf}
\caption{Integrated diagram of the Autonomy--Flow--Heat framework. Arrows denote
the causal dependencies proposed by the model. The structure mirrors an electrical
circuit with a driving potential and resistive dissipation, but the mapping is
analogical; the diagram is not intended to imply that organizational quantities
are measurable in physical units.}
\label{fig:model}
\end{figure}


\section{Decision, Orchestration, and Commit Architecture}
\label{sec:architecture}

Decision architecture is defined in this framework as the structural
configuration that allocates standing authority across an organizational
system: who may decide, within which domain, under which recognized mandate,
and with what escalation relationships. In the present model, decision
architecture primarily shapes legitimacy and the distribution of recognized
decision rights. It answers the question of \emph{who is structurally entitled
to decide in principle}.

Decision architecture alone, however, does not determine whether intent
survives translation into action. Organizations rarely act directly on
decision rights. They act through workflows. For this reason, the present
framework distinguishes decision architecture from \emph{orchestration
architecture}: the structural configuration through which decisions are
decomposed, routed, sequenced, escalated, reconciled, and translated into
coordinated action across human actors, software systems, and AI agents.
Orchestration architecture is the primary design lever for
$\text{Resistance}_{\text{Complexity}}$ and the principal intermediate layer
at which coherence is either preserved or lost.

A further distinction is required. Even when decision rights are clear and
workflow orchestration remains coherent, a specific act may still no longer be
valid when execution occurs. For that reason, decision architecture and
orchestration architecture must both be distinguished from \emph{commit
architecture}: the structural mechanism through which admissibility is
resolved at the point where action becomes materially, financially,
technically, or institutionally irreversible.

We denote execution-time admissibility for a specific act $a$ at time $t$ as
\[
\mathcal{A}(a,t) \in \{0,1\},
\]
where $\mathcal{A}(a,t)=1$ indicates that the act remains valid under current
conditions at the commit point, and $\mathcal{A}(a,t)=0$ indicates that it
does not. This notation is intentionally minimal. The present paper does not
yet specify the full functional form of admissibility, but introduces it to
mark the conceptual distinction between durable standing legitimacy and
execution-time validity.

\subsection{Relationship to Governance, Management, and Control Systems}
\label{sec:governance}

The three architectural layers introduced here operate at a different level of
abstraction from governance, management, and control systems. Governance
specifies who has ultimate authority. Management specifies how that authority
is exercised in planning, coordination, and supervision. Control systems
specify how deviations from desired states are detected and corrected. The
present framework introduces a more structural decomposition.

Decision architecture specifies how standing authority is distributed across
roles and domains. Orchestration architecture specifies how those decisions
are transformed into routable, sequenced, and enforceable action paths.
Commit architecture specifies how the system determines whether a specific act
may still proceed when execution becomes irreversible. Together these layers
define not merely who governs and who acts, but how intent travels from
recognized authority into operational reality and whether that action remains
valid when it commits.

This construct is conceptually related to Beer's viable systems model
\citep{beer1972brain}, in which system-level cohesion depends on the
regulatory structure of the metasystem. It is also related to Williamson's
treatment of organizational form as a determinant of transaction costs
\citep{williamson1981economics}. The present model extends these frameworks by
separating the distribution of authority from the preservation of intent
through orchestration and from the execution-time resolution of admissibility.
Whether this layered decomposition adds predictive value over more traditional
governance and coordination-cost accounts is an open empirical question.

\subsection{Workflow and Orchestration as the First Instability Layer}
\label{sec:orchestration_instability}

The framework proposes that in many complex organizations the first
operational instability appears in the workflow or orchestration layer rather
than in execution itself. The shop floor, customer interface, service outcome,
or compliance event typically reveals the symptom, but not the origin.
Similarly, data discrepancies, KPI conflicts, or reporting anomalies often
reflect already-distorted workflow logic rather than causing it. The
orchestration layer is therefore analytically privileged: it is where
strategic and managerial intent is first converted into routable, sequenced,
and enforceable action logic, and where small reinterpretations can become
systemically durable.

This claim has a specific implication. A workflow may remain locally coherent
at each step while globally encoding a distorted version of the original
decision. Handoffs may be clear, approvals correctly assigned, and local
metrics satisfied, yet the overall sequence may pursue a slightly altered
objective because the decomposition, routing, sequencing, escalation logic, or
reconciliation structure has already shifted the effective meaning of the
decision. In such cases, downstream execution is not malfunctioning in the
ordinary sense; it is functioning faithfully relative to a workflow that has
already drifted.

The framework therefore treats orchestration architecture not as a minor
implementation detail but as the principal structural layer through which
intent survives or degrades. This is also the layer where human and machine
action most directly interact: workflow engines, task routing systems,
approval paths, user interfaces, exception handling, vendor handoffs, and
agentic AI systems all operate here. As a result, this layer is expected to
become increasingly decisive in AI-augmented organizations.

\subsection{Commit Architecture and the Point of Irreversible Action}
\label{sec:commit}

Commit architecture is defined here as the structural arrangement by which the
system determines whether a specific act remains admissible at the point where
execution becomes materially, financially, technically, or institutionally
irreversible. Typical commit points include code deployment, release of funds,
grant of access, publication of instructions, execution of a trade, issuance
of an approval that triggers downstream automation, or activation of an
AI-originated action in a live environment.

This distinction matters because upstream governance artifacts---policies, role
assignments, workflow approvals, and documentation---shape and constrain
action but do not by themselves constitute control in the strong sense. If the
validity of the act is not resolved where the system commits, those upstream
elements function as preparation rather than decisive control. The model
therefore treats commit architecture as a separate and necessary complement to
both decision architecture and orchestration architecture.

A robust commit architecture must do at least three things. First, it must
identify the commit point clearly enough that irreversible action is not
occurring invisibly elsewhere in the stack. Second, it must make the
admissibility conditions for that class of action explicit enough to be
evaluated under current conditions. Third, it must bind the result of that
evaluation directly to execution, so that inadmissible actions are blocked
rather than merely recorded, questioned, or retrospectively audited.

This concept is especially important in distributed and AI-augmented
environments, where the interval between preparation and execution may be
short, the number of candidate actions large, and the contextual conditions
governing admissibility subject to rapid change.

\subsection{Structural Sources of Resistance}
\label{sec:resistance_sources}

The following structural properties are proposed to increase Complexity
Resistance, with special emphasis on the orchestration layer as the place
where these burdens are most often operationalized:

\begin{itemize}
  \item \textbf{Authority ambiguity:} role conflict and boundary ambiguity
    are empirically associated with performance decrements and stress
    \citep{rizzo1970role}; the model interprets this as a resistance increase.
  \item \textbf{Responsibility overlap:} diffusion of accountability is
    associated with reduced individual effort \citep{latane1979many}; the
    model treats this as a coordination friction that frequently becomes
    operationally embedded in workflow handoffs and exception paths.
  \item \textbf{Information--authority separation:} when decision rights are
    held by agents who lack the specific knowledge required to exercise them,
    information must be transferred before action can proceed
    \citep{jensen1992specific}; this transfer cost enters as resistance and
    often appears concretely as orchestration burden.
  \item \textbf{Incentive misalignment:} divergent incentive structures among
    interacting agents generate negotiation costs at decision nodes
    \citep{holmstrom1991multitask}, often forcing workflow branches,
    informal workarounds, or repeated reconciliation.
  \item \textbf{Feedback latency:} delayed feedback is a documented source of
    oscillation and instability in decision systems \citep{sterman1989modeling};
    the model treats it as a resistance component that increases the
    probability of constraint violations and stale orchestration logic.
  \item \textbf{Orchestration distortion:} decomposition logic, routing rules,
    sequencing assumptions, escalation pathways, and handoff structures may
    encode a slightly altered version of the original decision. The present
    framework treats this as a primary source of coherence loss rather than as
    a downstream symptom.
\end{itemize}

\subsection{Principles of Low-Resistance and High-Integrity Architecture}
\label{sec:principles}

From the model's assumptions, the following design principles reduce
Complexity Resistance and are therefore predicted to improve flow stability,
reduce heat, and preserve intent more reliably:

\begin{enumerate}
  \item Assigning clear, non-overlapping decision authority for each
    recognized decision class.
  \item Distinguishing explicitly between decision rights and the workflow
    structures through which those rights are operationalized.
  \item Collocating authority with the specific knowledge required to
    exercise it \citep{jensen1992specific}.
  \item Designing orchestration logic---decomposition, routing, sequencing,
    escalation, and reconciliation---so that the original decision can be
    translated into action without silent reinterpretation.
  \item Designing feedback mechanisms with sufficiently low latency to
    support the required throughput rate and to detect workflow drift before
    it becomes embedded in downstream data and execution.
  \item Resolving constraint conflicts at the architectural design level, so
    that individual decision nodes and workflows receive consistent,
    satisfiable constraint sets.
  \item Limiting the number of simultaneous incoming constraints at any node
    to within its estimated processing capacity.
  \item Making the point of irreversible commitment explicit and binding
    admissibility checks directly to that point, so that validity is resolved
    where execution occurs rather than assumed from earlier preparation.
\end{enumerate}

These are predictions of the model. The general principles are consistent with
established design guidance \citep{mintzberg1979structuring,
williamson1981economics}, but the specific prediction that violating them
generates superlinear heat and early orchestration instability---rather than
merely linear coordination cost---is a novel claim that requires empirical
support.


\section{Capability and Risk}
\label{sec:capability}

\[
\text{Capability} = \text{Capability}_{\max} - \text{Risk}
\]

Capability denotes the effective ability of an agent to execute actions
reliably under prevailing conditions. The framing as a situationally bounded
quantity reflects the well-established finding that performance depends not
only on skill but on the demands imposed by the environment: the same agent
may be highly capable under routine conditions and substantially impaired
under conditions of high uncertainty, cognitive load, or consequence severity.

$\text{Capability}_{\max}$ represents the agent's full repertoire of skills,
knowledge, and resources under low-demand conditions. Risk in this model
is not a property of outcomes per se but a demand on cognitive and
attentional resources. This framing draws on Kahneman's model of
attentional capacity \citep{kahneman1973attention}, which establishes that
higher-demand tasks reduce the resources available for concurrent activities,
but extends it to the organizational level in a way that source does not
directly imply.

The following factors are proposed as determinants of the Risk demand variable:

\begin{itemize}
  \item \textbf{Environmental uncertainty:} unpredictability increases
    monitoring demands and error-detection requirements \citep{perrow1984normal}.
  \item \textbf{Task complexity:} higher intrinsic task complexity
    consumes working memory capacity \citep{sweller1994cognitive}, leaving
    less available for productive execution.
  \item \textbf{Consequence severity:} high-stakes decisions direct
    attention toward loss-avoidance \citep{kahneman1979prospect}, reducing
    resources available for efficient action generation.
  \item \textbf{Knowledge gaps:} operating under conditions of bounded
    rationality \citep{simon1972theories} increases the cognitive cost of
    each decision, which the model treats as a reduction in net Capability.
\end{itemize}

As with the Legitimacy equation, the linear subtraction form is a modeling
simplification. The functional relationship between risk factors and
effective capability is almost certainly nonlinear in practice, and the
specific choice of linear subtraction is made for analytical tractability.
The parallel algebraic structure of the Legitimacy and Capability equations
is also a deliberate formal symmetry: both components of Autonomy are
modeled as bounded resources subject to environmental depletion. Whether
this symmetry reflects a deeper empirical regularity or is merely a
convenient formal property of the model is an open question.

\section{Action Flow}
\label{sec:flow}

\[
F = \frac{\text{Autonomy}}{\text{Resistance}_{\text{Complexity}}}
\]

Action flow $F$ represents the rate at which decisions and coordinated actions
are generated and completed per unit time within the organizational system.
The ratio form is directly inspired by Ohm's law in circuit theory
\citep{feynman1963lectures}, where current equals voltage divided by resistance.
This is an analogy, not a derivation: there is no prior theoretical or empirical
result that directly implies this functional form for organizational systems.
The ratio is proposed because it generates clear, testable predictions about
the effects of autonomy expansion and resistance reduction as independent
and interacting interventions.

Action flow should not be equated with valid commitment. A decision, instruction, or prepared action may move efficiently through an organizational structure without yet being executable under current conditions. For this reason, the present framework treats flow as the movement of coordinated action potential through the system, while reserving the question of whether a specific act may actually proceed for later treatment under the concept of admissibility at commitment.

$\text{Resistance}_{\text{Complexity}}$ subsumes all structural and
informational factors that impede the throughput of decisions and actions.
The following categories of resistance are proposed, each supported by
evidence that they increase coordination cost or reduce decision throughput:

\begin{itemize}
  \item \textbf{Coordination requirements:} the need to achieve alignment
    among multiple actors before action can proceed has been established
    as a systematic source of organizational cost \citep{malone1994interdisciplinary}.
  \item \textbf{Information overload:} processing demand in excess of
    working memory capacity degrades decision quality and speed
    \citep{sweller1988cognitive}.
  \item \textbf{Constraint conflict:} simultaneous obligations that cannot
    all be satisfied by a single decision force negotiation or escalation,
    both of which consume resources that would otherwise support throughput.
  \item \textbf{Structural opacity:} ambiguous authority maps and missing
    escalation paths are documented sources of coordination failure
    \citep{mintzberg1979structuring}.
\end{itemize}

The ratio form implies that flow can be increased either by increasing
Autonomy or by decreasing Resistance. These are architecturally distinct
interventions with different cost profiles. However, the critical implication
of the ratio form---in combination with the heat equation---is that
high Autonomy under high Resistance does not produce high stable flow;
it produces high organizational heat (Section~\ref{sec:heat}).
Whether this prediction holds quantitatively, as opposed to qualitatively,
is an empirical question the model cannot answer from its assumptions alone.

\section{Organizational Heat}
\label{sec:heat}

\[
H = F^2 \times \text{Resistance}
\]

Organizational heat $H$ is the central novel construct of this framework.
It is defined as the rate of unproductive dissipation within the
organizational system---manifesting as stress, interpersonal conflict,
rework, escalation, and decision fatigue. The term and its formal expression
are direct analogies to Joule heating in electrical systems
\citep{feynman1963lectures}, where power dissipation equals the square
of current times resistance. This is a conceptual innovation of the present
model and is not a standard construct in any of the literatures cited for
motivation.

The quadratic dependence on flow is the most consequential structural
feature of the model. Under this formulation, doubling flow while holding
resistance constant quadruples heat. This non-linearity generates
a specific prediction: systems that operate stably under moderate action
rates may exhibit abrupt, disproportionate degradation under higher loads.
This qualitative pattern is documented across several literatures.
The Yerkes--Dodson relationship between arousal and performance
\citep{yerkes1908relation} establishes an inverted-U pattern in which
high arousal impairs performance---consistent with the prediction that high
flow generates costs---but does not imply a quadratic relationship specifically.
Perrow's analysis of tightly coupled systems \citep{perrow1984normal}
documents how small perturbations can cascade to catastrophic failure,
which is consistent with superlinear heat accumulation but is not formalized
in those terms. The specific quadratic form of the heat equation is a modeling
assumption that warrants direct empirical investigation.

This formulation captures dissipation generated by friction in the movement and reconciliation of coordinated action. It does not by itself capture a distinct class of failure in which an action reaches the point of execution efficiently but should not occur because its authorizing conditions are no longer satisfied. Such failures are not best understood as high heat alone, but as failures of admissibility resolution at commitment. The two may interact: high heat increases the likelihood of stale, bypassed, or weakly validated commitments, but the concepts should remain analytically distinct.

Substituting the flow equation yields the alternative expression:

\[
H = \frac{\text{Autonomy}^2}{\text{Resistance}_{\text{Complexity}}}
\]

This form reveals an important design asymmetry: under the model's
assumptions, increasing Autonomy without reducing Resistance raises heat
superlinearly, while reducing Resistance (holding Autonomy constant) reduces
heat. The practical implication is that resistance reduction is a more
leverage-efficient intervention than autonomy expansion under high-complexity
conditions. This is a prediction of the model, not an established finding.

When heat is not dissipated through structural adjustment, the model
predicts accumulation as latent instability. The psychological correlates
proposed---burnout \citep{maslach1981measurement}, trust erosion
\citep{mayer1995integrative}, and risk aversion---are all empirically
documented outcomes of sustained organizational overload. The model proposes
that they share a common structural cause expressible through the heat
equation; whether that causal unification is valid is an empirical question.

\section{Kirchhoff Constraints in Decision Networks}
\label{sec:kirchhoff}

The scalar relations of Sections~\ref{sec:flow}--\ref{sec:heat} describe
individual agents or units in isolation. Real organizations are networked:
decisions depend on prior decisions, authorities overlap, and information
flows along paths. To model these network effects, we introduce
Kirchhoff-like conservation constraints by direct analogy to Kirchhoff's
current and voltage laws in electrical network theory \citep{feynman1963lectures}.

These constraints are the authors' extension of the flow analogy to networked
systems. They are not derived from organizational theory but are proposed
as useful formal tools for identifying structural sources of coordination
failure. Their value depends on whether the predictions they generate---
developed in Section~\ref{sec:predictions}---are empirically supported.

\subsection{Node Conservation: Flow Balance}

We propose that at any decision node $n$ in the organizational network,
a conservation-like constraint should hold between incoming and outgoing
decision flow:

\[
\sum_{i} F_{\text{in},i} = \sum_{j} F_{\text{out},j}
\]

A decision node represents any locus where multiple information streams,
authorities, or constraints must be reconciled before an output decision
can be produced. Common examples include a project manager integrating
inputs from engineering, finance, and compliance before committing to a
delivery plan, or a cross-functional committee reconciling competing
resource claims.

The proposed conservation principle states that if incoming flows
consistently exceed the node's processing capacity, the excess accumulates
as a backlog of unresolved decisions. This backlog is itself a source of
increased Resistance for subsequent flow, creating a positive feedback:
backlog increases resistance, which reduces effective throughput, which may
increase heat as upstream agents attempt to force decisions through the
congested node. This feedback structure is qualitatively consistent with
documented patterns of decision escalation and overload in management
research \citep{sterman1989modeling}, but the formal correspondence has
not been empirically established.

Decision nodes must reconcile all incoming constraints simultaneously,
which may include strategic goals, resource limits, risk and compliance
policies, stakeholder demands, and temporal dependencies. When the
incoming constraint set is overdetermined---when no single decision can
simultaneously satisfy all constraints---the node cannot produce a valid
outgoing flow. The model characterizes this as a node-level conservation
violation. In practice, such violations manifest as decision deadlock or
escalation; the model predicts that nodes with chronically overdetermined
constraint sets will be persistent heat sources. This prediction is testable
and is stated formally in Section~\ref{sec:predictions}.

\subsection{Loop Constraints: Potential Balance}

For closed loops within the decision network, we propose the following
constraint by analogy to Kirchhoff's voltage law:

\[
\sum_{\text{loop}} \text{Autonomy}_k = \sum_{\text{loop}} F_k \times \text{Resistance}_k
\]

This constraint states that the total driving potential invested in a
decision cycle must be fully accounted for by productive work done against
resistance within that cycle. When the balance is violated---when more
potential is invested than the cycle can convert to output---the residual
accumulates as heat or forces revision of the assumptions underlying the cycle.

Organizational loops subject to potential imbalance include:

\begin{itemize}
  \item \textbf{Strategy--execution--evaluation--strategy cycles:}
    where strategic ambition is not matched by the capability and resistance
    profile of the execution layer, producing persistent performance gaps
    that the feedback mechanism cannot correct \citep{sterman1989modeling}.
  \item \textbf{Policy--implementation--reporting--policy cycles:}
    where compliance reporting is structurally decoupled from implementation
    realities, producing cycles that systematically miscalibrate the
    policy layer.
  \item \textbf{Objective--action--feedback--objective cycles:}
    where feedback latency or distortion prevents accurate potential signals
    from reaching the goal-setting process.
\end{itemize}

\subsection{Coordination Failures as Conservation Violations}

The framework proposes that many coordination failures can be formally characterized as violations of node or loop conservation. A unit that generates more outgoing decisions than its incoming standing authority and constraints can support is operating with a legitimacy or coherence deficit. A unit that absorbs incoming decision flow without producing proportionate outputs is generating a resistance backlog. Both conditions degrade systemic flow efficiency and increase total organizational heat. However, satisfaction of these conservation-like constraints is not sufficient for full control, because a network may remain internally balanced while still committing an inadmissible act if validity is not resolved at the point where execution becomes irreversible.

This characterization has diagnostic value: it suggests that persistent
coordination failures can be localized to specific nodes or loops where
conservation is violated, and that structural interventions targeting
those points should be more effective than diffuse cultural or motivational
interventions. Whether this diagnostic strategy outperforms alternatives
in practice is an empirical question that the present framework cannot
resolve on theoretical grounds alone.

\section{AI-Native Organizations}
\label{sec:ai}

\subsection{How AI Augments Potential Flow}

AI-augmented organizations---systems in which artificial intelligence performs
a substantial fraction of information processing, decision-support, or
execution functions---alter the model's parameters in ways that are
qualitatively distinctive. There is empirical evidence that AI tools extend
the range of tasks that can be executed reliably \citep{agrawal2018prediction}
and reduce the time required per decision, which in the model's terms
increases Autonomy at the agent or unit level by reducing Risk demand
on human Capability. Automation similarly reduces variance in routine
execution, decreasing the uncertainty component of Risk.

Under the model's flow equation $F = \text{Autonomy} / \text{Resistance}$,
these effects produce a substantial increase in potential action flow.
Whether that potential is realized as productive output or dissipated
as heat depends on whether Resistance is adjusted proportionately.
The model's prediction---and this is explicitly a model-derived prediction
rather than an observed regularity---is that when AI increases Autonomy
without a proportionate reduction in Resistance, the heat equation
$H = F^2 \times \text{Resistance}$ predicts superlinear heat accumulation.

\subsection{Why Heat May Increase Faster Than Flow}

The model predicts that AI-augmented systems face an asymmetric risk:
Autonomy-increasing mechanisms and Resistance-reducing mechanisms
operate independently and are not automatically co-produced. Several
properties of AI systems are proposed to generate resistance of their own:

\begin{itemize}
  \item \textbf{Objective misalignment:} AI systems optimizing proxy
    metrics can diverge from organizational intent, creating constraint
    conflicts at integration nodes \citep{russell2019human}. The model
    treats each such conflict as a Kirchhoff node violation.
  \item \textbf{Metric proliferation:} multiple AI subsystems generating
    competing performance signals increase constraint overdetermination
    at human oversight nodes.
  \item \textbf{Opacity, legitimacy loss, and weak admissibility resolution:} non-interpretable AI decisions may reduce the social legitimacy of AI-generated standing within the human organizational system, while also making it more difficult to determine whether a specific AI-originated act remains admissible under current conditions. The model formalizes the first as an increase in Insignificance for AI-origin flows and treats the second as a commit-architecture problem.
  \item \textbf{Pace asymmetry:} AI subsystems may generate decisions at
    rates that exceed the processing capacity of human governance nodes,
    producing backlogs that the model identifies as conservation violations.
  \item \textbf{Cascading dependencies:} tightly coupled AI pipelines
    increase the topological complexity of the decision network, multiplying
    the number of Kirchhoff constraints that must be simultaneously satisfied.
\end{itemize}

These effects are proposed as mechanisms by which AI integration could
generate heat faster than flow. Whether they do so in practice, and
whether the magnitude follows the quadratic prediction, are empirical
questions that organizational research on AI-augmented work could
in principle address.

\subsection{Architectural Requirements for AI-Augmented Stability}

The model generates the following specific structural predictions for
AI-augmented organizations:

\begin{itemize}
  \item Each AI agent requires a clearly specified authority domain, understood as the class of actions or recommendations for which it may generate prepared outputs under normal conditions.
  \item Irreversible AI-originated actions must be subject to execution-time admissibility checks bound directly to the commit point, rather than inferred solely from prior approvals, static rules, or durable role assignments.
  \item The objective functions of interacting AI agents must be mutually consistent to avoid loop-level potential imbalance and locally valid but globally misaligned action.
  \item Human oversight nodes must have sufficient processing capacity to handle the flow generated by AI subsystems, or AI decision rates must be deliberately throttled to match human governance bandwidth.
  \item Feedback latency must be low enough to detect objective drift, contextual change, and constraint conflict before they accumulate as unresolved backlogs or invalid commitments.
\end{itemize}

These are architectural prescriptions derived from the model's assumptions.
They are consistent with the governance recommendations of \citet{russell2019human}
and with emerging frameworks for AI safety and controllability, but their
specific derivation from the heat equation is novel to this framework.

In AI-augmented settings, the distinction between standing legitimacy and admissibility becomes operationally decisive. A model may be correctly assigned to a decision domain, yet any particular act it advances toward execution may still be invalid when current constraints, permissions, exposures, or environmental conditions are taken into account. The control problem therefore cannot be solved by authority specification alone; it requires commit architecture capable of recomputing or re-evaluating admissibility at execution time.

\section{The Coordination Economy}
\label{sec:coordination}

\subsection{Coordination Cost as the Binding Constraint}

The term \textit{coordination economy} is used here to describe a class of
productive environment in which the technical means of production are not the
primary constraint on value creation, and the central bottleneck lies instead
in preserving coherent and admissible action across distributed agents. This
environment is characteristic of knowledge-intensive industries, platform
ecosystems, regulated organizations, and multi-agent AI systems, where many
actors can generate analysis, recommendations, and even execution-ready
actions at low marginal cost, but where the difficulty of maintaining aligned
direction and valid commitment grows with system scale, interdependence, and
speed \citep{malone1994interdisciplinary, parker2016platform,
puranam2012coordination}.

In this sense, coordination cost should not be understood narrowly as the cost
of communication or synchronization. In high-complexity environments, the
deeper problem is that action must satisfy two distinct structural conditions
at once. First, intent must survive translation into workflow. Strategic and
managerial decisions must be decomposed, routed, sequenced, escalated, and
reconciled without silently altering what the organization is actually trying
to do. Second, a specific act must remain admissible at the moment of
execution: the authority, context, and constraint set under which the action
becomes irreversible must still render that action valid under current
conditions.

The central bottleneck in a coordination economy is therefore not merely
moving information or obtaining approvals. It is preserving intent through
orchestration while still resolving admissibility at commitment. Where this
capability is weak, organizations can generate abundant decisions, analyses,
and workflow activity while still failing to realize valid coordinated action.

\subsection{From Standing Authority to Orchestrated Action}

The earlier sections model legitimacy as a structural component of autonomy: a
relatively durable property of recognized standing within an organizational
decision system. That concept remains necessary, but it is not sufficient.
Possessing standing authority does not ensure that the authority will be
translated into action faithfully. Between recognized decision rights and
irreversible execution lies an intermediate layer: workflow or orchestration
architecture.

This layer determines how decisions are operationalized across handoffs,
systems, teams, and increasingly AI agents. It determines how a decision is
broken into sub-decisions, how tasks are routed, how dependencies are ordered,
how exceptions are escalated, and how conflicting constraints are reconciled.
In doing so, it becomes the principal site where coherence is either preserved
or lost. The present framework therefore proposes that data problems and
execution breakdowns are often downstream observability surfaces of
orchestration failure rather than the primary origin of drift.

The coordination economy is therefore not only an economy of alignment across
distributed actors; it is an economy of preserving intent as decisions are
translated into workflows. A system with clear authority but weak
orchestration integrity may remain procedurally orderly while advancing a
distorted version of the original objective. Conversely, a system with strong
orchestration but weak standing legitimacy may route work efficiently while
lacking recognized mandate or acceptance. Durable coordinated performance
requires both.

\subsection{Why Flow, Resistance, and Heat Matter at Scale}

The flow--resistance--heat framework developed in Sections~\ref{sec:flow}
and~\ref{sec:heat} remains central to this analysis. It models the
organizational cost of moving decisions and actions through complex
structures. In low-complexity, low-throughput environments, the heat
generated by suboptimal architecture may remain modest. As throughput,
interdependence, and coupling increase, however, the architecture's
resistance profile becomes increasingly consequential. Under the model's
assumptions, systems with higher resistance generate more dissipation,
greater backlog, and more frequent escalation for the same nominal level of
action flow. The result is that architectural advantages compound with
scale.

In coordination-economy environments, however, throughput alone is not the
relevant measure of success. A system may exhibit high flow in the sense of
rapid decision preparation, routing, and authorization, while still failing
to preserve the original intent through the orchestration layer. In such
cases, rework, escalation, conflict, and fatigue may arise not only from high
volume under resistance, but from repeated attempts to force throughput
through workflows that already encode a drifted objective. Heat therefore
captures not only overload but also dissipation generated by compensating for
distorted coordination logic.

For this reason, the present framework distinguishes analytically between
\textit{coordination flow} and \textit{commitment}. Coordination flow refers
to the movement of intent, information, authority claims, and preparatory
decisions through the organizational network. Commitment refers to the moment
at which action becomes materially or institutionally irreversible---for
example, deployment of code, release of funds, grant of access, issuance of
an instruction, or execution of an AI-originated action in a live
environment.

This distinction matters because the structural conditions required for high
flow are not identical to those required for valid commitment. A system may
be optimized for throughput while still committing inadmissible actions if
the final execution point does not resolve whether the action is still valid
under current conditions. In such cases, earlier governance, documentation,
and workflow steps function only as preparation. They may shape the action,
but they do not control it in the strong sense unless admissibility is
resolved where the system commits.

The coordination economy therefore operates under a dual structural burden.
Organizations must reduce resistance sufficiently to sustain coherent flow
through the orchestration layer, but they must also ensure that irreversible
action is conditioned on execution-time validity. The first burden is
captured by the resistance and heat relations introduced earlier. The second
burden introduces an additional architectural requirement: the system must not
merely transport authority claims; it must determine whether action is still
admissible when execution occurs.

\subsection{Preparation, Orchestration, Commitment, and Control}

A central implication of the present framework is that preparation should not
be conflated with control, and that orchestration should not be mistaken for a
purely technical implementation layer. Organizations invest heavily in
preparatory structures: governance committees, policy documents, approval
chains, workflow steps, model documentation, and audit trails. These may shape
the generation and routing of action significantly. But unless they preserve
intent through orchestration and bind the point at which irreversible action
occurs, they do not control execution in the strongest sense.

This distinction matters because many organizational pathologies arise
precisely when preparation is mistaken for control. A system may appear
well-governed because decisions are documented, authority domains are
assigned, workflows are formalized, and approvals are captured, while still
allowing either of two structural failures. First, the workflow may have
already encoded a distorted version of the original decision. Second, an
inadmissible act may still commit because no mechanism resolves current
validity at execution time. In such cases, the organization has designed the
preparation path and perhaps even the process map, but not the full control
path.

The coordination economy therefore depends on three related but non-identical
architectures. The first is the \emph{decision architecture} that assigns
standing authority. The second is the \emph{orchestration architecture} that
preserves coherence, reduces resistance, and translates decisions into
coordinated action paths. The third is the \emph{commit architecture} that
determines whether the act remains admissible under current conditions at the
moment the system commits. Durable advantage requires all three.

\subsection{Two Distinct Failure Modes: Orchestration Failure and Commit Failure}

This refinement clarifies that large-scale coordination failure can arise
through two analytically distinct mechanisms, which may occur independently or
jointly.

First, systems may fail through \textit{orchestration failure}. In this case,
the original strategic or operational intent is progressively reinterpreted as
it moves through decomposition logic, routing rules, handoffs, sequencing
structures, escalation paths, vendor interfaces, software systems, and
automation pipelines. Each local translation may appear reasonable, and each
step may remain procedurally valid, yet the action ultimately committed is no
longer aligned with the intent the organization sought to realize. The failure
here is not necessarily stale authority but structural drift: the system
preserves permissions, process conformance, or local optimization while losing
global directional integrity. The first practical instability often appears in
the orchestration layer, while execution and data reveal the consequences only
later.

Second, systems may fail through \textit{commit failure}. In this case, the
intended action may remain coherent and the authority under which it was
prepared may initially have been valid, but by the moment of execution the
conditions that justified the action no longer hold. The context has changed,
the authorization has expired, the risk state has shifted, or a higher-order
constraint now blocks execution. The failure here is temporal rather than
translational: the problem is not that the action was distorted in transit,
but that it was not revalidated where it became irreversible.

These failure modes may occur independently or jointly. A system may preserve
orchestration integrity yet execute stale authority. A system may preserve
admissibility at execution time yet coordinate around a distorted objective.
Or it may do both: drift structurally while also failing to validate
commitment. The coordination economy is therefore not reducible to either
throughput optimization or compliance checking alone. Its central problem is
the joint preservation of coherent intent through orchestration and
admissible action at commitment under distributed conditions.

A system may therefore be structurally coherent yet inadmissible at
commitment, or admissible at commitment yet globally incoherent with respect
to the organization's broader intent.

\subsection{Why AI Increases the Importance of Orchestration and Commit Architecture}

These issues become particularly acute in AI-augmented organizations.
Artificial intelligence can increase the effective rate at which options are
generated, evaluated, and advanced toward execution. As argued in
Section~\ref{sec:ai}, this can increase potential flow substantially. But AI
also compresses the interval between preparation and commitment, scales the
number of actions under consideration, and increases the number of translation
points through which organizational intent must pass before action occurs.

As a result, AI systems intensify both major coordination burdens. They may
increase orchestration risk by introducing additional translation layers
between organizational intent and operational action, especially when
objective functions, proxy metrics, workflow engines, or learned policies
substitute for direct human judgment. They also increase admissibility risk by
enabling action at a pace that exceeds the organization's capacity to
determine whether authority and context remain valid at the point of
execution. Under these conditions, a governance arrangement that remains
external to the orchestration and execution layers may shape behavior but
cannot reliably control it.

The model therefore predicts that in AI-native coordination economies,
competitive advantage will not accrue simply to organizations that generate
more decisions, automate more tasks, or distribute more autonomy. It will
accrue to organizations that build architectures capable of doing three
things simultaneously:

\begin{enumerate}
  \item preserving the coherence of intent through workflow and orchestration,
  \item minimizing structural resistance so that flow remains stable rather
    than heat-generating, and
  \item resolving admissibility at the point of irreversible execution.
\end{enumerate}

\subsection{Implications for Organizational Advantage at Scale}

Under these conditions, the central unit of competition shifts. It is no
longer sufficient to compare organizations by production efficiency,
information access, or local decision quality alone. The relevant question
becomes whether the organization can preserve valid coordinated action under
conditions of high throughput, distributed interdependence, and changing
context.

The model therefore predicts that organizations with lower structural
resistance, stronger orchestration integrity, and stronger execution-time
validity mechanisms will realize compounding advantages at scale.
Specifically, it predicts that:

\begin{itemize}
  \item organizations that expand autonomy without reducing resistance will
    generate heat rather than proportional value;
  \item organizations that formalize decision rights without preserving
    intent through orchestration will exhibit orderly but drifted execution;
  \item organizations that improve throughput without embedding
    execution-time admissibility checks will increase the risk of invalid
    commitments;
  \item organizations that maintain local validity while allowing intent to
    fragment through workflow layers will exhibit globally misaligned but
    procedurally correct execution;
  \item organizations that jointly optimize decision architecture,
    orchestration architecture, resistance reduction, and commit
    architecture will be able to sustain higher effective throughput before
    encountering instability thresholds.
\end{itemize}

The coordination economy, as defined here, therefore refers not merely to an
environment in which coordination costs are high. It refers to an environment
in which value creation increasingly depends on an organization's ability to
preserve coherent intent through orchestration while also resolving
admissibility at the point of irreversible execution. In such environments,
resistance management, orchestration integrity, network conservation, and
commit-time admissibility become coequal determinants of organizational
viability.


\section{Limitations of the Model}
\label{sec:limitations}

The following limitations are inherent to the framework as currently formulated
and should be addressed in any research program designed to evaluate its claims.

\begin{enumerate}
  \item \textbf{Unvalidated measurement:} The model's constructs---Legitimacy,
    Capability, Insignificance, Risk, and Resistance---are defined
    conceptually but do not yet have validated psychometric or organizational
    measurement instruments. The algebraic precision of the expressions
    is not currently matched by corresponding precision in operationalization.
    Construct validation studies are a necessary precondition for empirical
    testing of any of the predictions in Section~\ref{sec:predictions}.

  \item \textbf{Assumed functional forms:} The specific functional
    forms---multiplicative for Autonomy, linear subtraction for Legitimacy
    and Capability, ratio for Flow, and quadratic for Heat---are modeling
    choices made for analytical tractability and conceptual coherence.
    Empirical relationships among the underlying variables may be nonlinear,
    discontinuous, threshold-dependent, or interaction-dependent in ways
    that these forms do not capture. The model makes no claim that these
    are the uniquely correct functional forms.

  \item \textbf{Static structure:} The framework describes system state
    at a point in time. Organizational systems are dynamic: resistance
    changes as architecture adapts, legitimacy changes as relationships
    develop, and capability changes through learning and attrition.
    The model does not currently specify differential equations or
    feedback dynamics governing these changes.

  \item \textbf{Interpersonal and cultural factors:} The framework models
    organizations as flow networks and does not represent interpersonal
    relationships, organizational culture, or historical path dependence---
    all of which are empirically significant \citep{schein1985organizational}.
    The model's resistance variable is intended to absorb some cultural
    effects, but this substitution is not theoretically grounded.

  \item \textbf{Kirchhoff analogy limitations:} Kirchhoff's laws are
    formally valid for linear, time-invariant networks. Organizational
    decision networks are nonlinear, time-varying, and subject to
    exogenous shocks. The conservation constraints introduced in
    Section~\ref{sec:kirchhoff} should be understood as qualitative
    structural hypotheses rather than precise quantitative laws.

  \item \textbf{Endogeneity:} The model specifies a causal structure
    in which architecture determines resistance, which determines flow
    and heat. In practice, heat outcomes (conflict, burnout, turnover)
    alter architectural properties, creating feedback loops not currently
    represented in the model. Causal identification in any empirical
    test will need to address this endogeneity through longitudinal
    or experimental designs.

  \item \textbf{Admissibility not yet formalized algebraically:} While the revised framework distinguishes standing legitimacy from execution-time admissibility, it does not yet provide a full algebraic treatment of admissibility comparable to the treatments of Autonomy, Flow, and Heat. The current paper therefore identifies admissibility as a necessary construct and commit architecture as a necessary design layer, but stops short of specifying a complete functional form for how contextual change, temporal decay, and commit-point validation interact.

  \item \textbf{The ``heat'' metaphor and reification risk:} The term
    organizational heat carries strong inferential connotations from
    physics that may not transfer reliably to organizational contexts.
    Readers and researchers should treat the term as a label for a
    proposed latent variable---unproductive dissipation under structural
    load---rather than as a claim that organizations literally obey
    thermodynamic laws.
\end{enumerate}

\section{Conclusion}

This paper has proposed a structural model of organizational dynamics
grounded in physical flow analogies and positioned within established
theoretical traditions in psychology, systems theory, and organization
design. The core framework formalizes Autonomy as the product of Legitimacy
and Capability, Action Flow as the ratio of Autonomy to Complexity
Resistance, and Organizational Heat as a quadratic function of Flow and
Resistance. These are modeling hypotheses---deliberately stated in formal
terms to enable testable predictions---not summaries of existing findings.

The framework is motivated by, but distinct from, its theoretical antecedents.
Self-determination theory, cognitive load theory, cybernetics, normal accident
theory, and coordination economics each supply evidence for one or more of
the model's conceptual components; none of them directly implies the specific
algebraic forms proposed here.

The model's primary theoretical contributions are fourfold. First, it proposes a unified formal language for phenomena---learned helplessness, burnout, coordination failure, escalation overload---that are typically treated in separate literatures. Second, it applies Kirchhoff-like conservation constraints to decision networks, generating structural predictions about where coordination failures will occur and why they are stable. Third, it distinguishes standing legitimacy from execution-time admissibility, thereby separating the problem of coordinated preparation from the problem of valid commitment. Fourth, it applies the flow--resistance--heat framework to AI-augmented organizations, generating predictions about the conditions under which AI-driven capability expansion produces organizational instability or invalid commitment rather than productivity.

The testable predictions developed in Section~\ref{sec:predictions} provide a concrete research agenda. Testing them would require both construct validation---developing reliable measures for the model's latent variables---and structural equation modeling, longitudinal designs, or experimental methods capable of separating coherence failure from admissibility failure and estimating the functional forms proposed. Until such tests are conducted, the framework should be understood as a structured theoretical proposal whose value lies in the precision and falsifiability of its claims, not in their prior empirical confirmation.

\bibliographystyle{apalike}
\begin{thebibliography}{99}

\bibitem[Agrawal et al., 2018]{agrawal2018prediction}
Agrawal, A., Gans, J., \& Goldfarb, A. (2018).
\textit{Prediction Machines: The Simple Economics of Artificial Intelligence}.
Harvard Business Review Press.

\bibitem[Ashby, 1956]{ashby1956introduction}
Ashby, W. R. (1956).
\textit{An Introduction to Cybernetics}.
Chapman \& Hall.

\bibitem[Beer, 1972]{beer1972brain}
Beer, S. (1972).
\textit{Brain of the Firm}.
Allen Lane.

\bibitem[Bertalanffy, 1968]{bertalanffy1968general}
Bertalanffy, L. von (1968).
\textit{General System Theory}.
George Braziller.

\bibitem[Cialdini, 2001]{cialdini2001influence}
Cialdini, R. B. (2001).
\textit{Influence: Science and Practice} (4th ed.).
Allyn \& Bacon.

\bibitem[Csikszentmihalyi, 1990]{csikszentmihalyi1990flow}
Csikszentmihalyi, M. (1990).
\textit{Flow: The Psychology of Optimal Experience}.
Harper \& Row.

\bibitem[Deci \& Ryan, 1985]{deci1985intrinsic}
Deci, E. L., \& Ryan, R. M. (1985).
\textit{Intrinsic Motivation and Self-Determination in Human Behavior}.
Plenum.

\bibitem[Deci \& Ryan, 2000]{deci2000self}
Deci, E. L., \& Ryan, R. M. (2000).
The ``what'' and ``why'' of goal pursuits: Human needs and the self-determination
of behavior.
\textit{Psychological Inquiry}, 11(4), 227--268.

\bibitem[Demerouti et al., 2001]{demerouti2001job}
Demerouti, E., Bakker, A. B., Nachreiner, F., \& Schaufeli, W. B. (2001).
The job demands--resources model of burnout.
\textit{Journal of Applied Psychology}, 86(3), 499--512.

\bibitem[Eisenberger et al., 1986]{eisenberger1986perceived}
Eisenberger, R., Huntington, R., Hutchison, S., \& Sowa, D. (1986).
Perceived organizational support.
\textit{Journal of Applied Psychology}, 71(3), 500--507.

\bibitem[Feynman et al., 1963]{feynman1963lectures}
Feynman, R. P., Leighton, R. B., \& Sands, M. (1963).
\textit{The Feynman Lectures on Physics} (Vol.\ 2).
Addison-Wesley.

\bibitem[Hackman \& Oldham, 1980]{hackman1980work}
Hackman, J. R., \& Oldham, G. R. (1980).
\textit{Work Redesign}.
Addison-Wesley.

\bibitem[Holmstr\"om \& Milgrom, 1991]{holmstrom1991multitask}
Holmstr\"om, B., \& Milgrom, P. (1991).
Multitask principal-agent analyses.
\textit{Journal of Law, Economics, and Organization}, 7, 24--52.

\bibitem[Jensen \& Meckling, 1992]{jensen1992specific}
Jensen, M. C., \& Meckling, W. H. (1992).
Specific and general knowledge, and organizational structure.
In L. Werin \& H. Wijkander (Eds.),
\textit{Contract Economics} (pp.\ 251--274).
Blackwell.

\bibitem[Kahneman, 1973]{kahneman1973attention}
Kahneman, D. (1973).
\textit{Attention and Effort}.
Prentice-Hall.

\bibitem[Kahneman \& Tversky, 1979]{kahneman1979prospect}
Kahneman, D., \& Tversky, A. (1979).
Prospect theory: An analysis of decision under risk.
\textit{Econometrica}, 47(2), 263--291.

\bibitem[Latan\'e et al., 1979]{latane1979many}
Latan\'e, B., Williams, K., \& Harkins, S. (1979).
Many hands make light the work: The causes and consequences of social loafing.
\textit{Journal of Personality and Social Psychology}, 37(6), 822--832.

\bibitem[Lazarus \& Folkman, 1984]{lazarus1984stress}
Lazarus, R. S., \& Folkman, S. (1984).
\textit{Stress, Appraisal, and Coping}.
Springer.

\bibitem[Malone \& Crowston, 1994]{malone1994interdisciplinary}
Malone, T. W., \& Crowston, K. (1994).
The interdisciplinary study of coordination.
\textit{ACM Computing Surveys}, 26(1), 87--119.

\bibitem[Maslach \& Jackson, 1981]{maslach1981measurement}
Maslach, C., \& Jackson, S. E. (1981).
The measurement of experienced burnout.
\textit{Journal of Occupational Behavior}, 2(2), 99--113.

\bibitem[Mayer et al., 1995]{mayer1995integrative}
Mayer, R. C., Davis, J. H., \& Schoorman, F. D. (1995).
An integrative model of organizational trust.
\textit{Academy of Management Review}, 20(3), 709--734.

\bibitem[Milgrom \& Roberts, 1990]{milgrom1990rationalizability}
Milgrom, P., \& Roberts, J. (1990).
Rationalizability, learning, and equilibrium in games with strategic
complementarities.
\textit{Econometrica}, 58(6), 1255--1277.

\bibitem[Mintzberg, 1979]{mintzberg1979structuring}
Mintzberg, H. (1979).
\textit{The Structuring of Organizations}.
Prentice-Hall.

\bibitem[Parker et al., 2016]{parker2016platform}
Parker, G. G., Van Alstyne, M. W., \& Choudary, S. P. (2016).
\textit{Platform Revolution}.
W. W. Norton.

\bibitem[Perrow, 1984]{perrow1984normal}
Perrow, C. (1984).
\textit{Normal Accidents: Living with High-Risk Technologies}.
Basic Books.

\bibitem[Puranam et al., 2012]{puranam2012coordination}
Puranam, P., Raveendran, M., \& Knudsen, T. (2012).
Organization design: The epistemic interdependence perspective.
\textit{Academy of Management Review}, 37(3), 419--440.

\bibitem[Rizzo et al., 1970]{rizzo1970role}
Rizzo, J. R., House, R. J., \& Lirtzman, S. I. (1970).
Role conflict and ambiguity in complex organizations.
\textit{Administrative Science Quarterly}, 15(2), 150--163.

\bibitem[Russell, 2019]{russell2019human}
Russell, S. (2019).
\textit{Human Compatible: Artificial Intelligence and the Problem of Control}.
Viking.

\bibitem[Ryan \& Deci, 2000]{ryan2000self}
Ryan, R. M., \& Deci, E. L. (2000).
Self-determination theory and the facilitation of intrinsic motivation,
social development, and well-being.
\textit{American Psychologist}, 55(1), 68--78.

\bibitem[Schein, 1985]{schein1985organizational}
Schein, E. H. (1985).
\textit{Organizational Culture and Leadership}.
Jossey-Bass.

\bibitem[Seligman, 1975]{seligman1975helplessness}
Seligman, M. E. P. (1975).
\textit{Helplessness: On Depression, Development, and Death}.
Freeman.

\bibitem[Simon, 1972]{simon1972theories}
Simon, H. A. (1972).
Theories of bounded rationality.
In C. B. McGuire \& R. Radner (Eds.),
\textit{Decision and Organization} (pp.\ 161--176).
North-Holland.

\bibitem[Staw, 1981]{staw1981escalation}
Staw, B. M. (1981).
The escalation of commitment to a course of action.
\textit{Academy of Management Review}, 6(4), 577--587.

\bibitem[Sterman, 1989]{sterman1989modeling}
Sterman, J. D. (1989).
Modeling managerial behavior: Misperceptions of feedback in a dynamic
decision making experiment.
\textit{Management Science}, 35(3), 321--339.

\bibitem[Sweller, 1988]{sweller1988cognitive}
Sweller, J. (1988).
Cognitive load during problem solving: Effects on learning.
\textit{Cognitive Science}, 12(2), 257--285.

\bibitem[Sweller, 1994]{sweller1994cognitive}
Sweller, J. (1994).
Cognitive load theory, learning difficulty, and instructional design.
\textit{Learning and Instruction}, 4(4), 295--312.

\bibitem[Weick, 1995]{weick1995sensemaking}
Weick, K. E. (1995).
\textit{Sensemaking in Organizations}.
Sage.

\bibitem[Wiener, 1948]{wiener1948cybernetics}
Wiener, N. (1948).
\textit{Cybernetics: Or Control and Communication in the Animal and the Machine}.
MIT Press.

\bibitem[Williamson, 1981]{williamson1981economics}
Williamson, O. E. (1981).
The economics of organization: The transaction cost approach.
\textit{American Journal of Sociology}, 87(3), 548--577.

\bibitem[Yerkes \& Dodson, 1908]{yerkes1908relation}
Yerkes, R. M., \& Dodson, J. D. (1908).
The relation of strength of stimulus to rapidity of habit-formation.
\textit{Journal of Comparative Neurology and Psychology}, 18(5), 459--482.

\bibitem[Yukl, 2012]{yukl2012leadership}
Yukl, G. (2012).
\textit{Leadership in Organizations} (8th ed.).
Pearson.

\end{thebibliography}

\end{document}
